% !TEX program = lualatex
\documentclass[12pt,a4paper]{article}

% === Moderne LuaLaTeX-Pakete (statt veralteter) ===
% ENTFERNT: inputenc, fontenc, babel (nicht nötig/problematisch mit LuaLaTeX)
\usepackage{fontspec}                    % Moderne Fonts für LuaLaTeX
\usepackage{polyglossia}                 % Moderne Sprachunterstützung
\setmainlanguage[spelling=new]{german}   % Neue deutsche Rechtschreibung
\usepackage{unicode-math}                % Unicode-Mathematik

% === Layout und Listen ===
\usepackage{enumitem}
\usepackage{tabularray}                  % Moderne Tabellen (statt longtable/array)
\usepackage{hyperref}
\usepackage{geometry}
\geometry{margin=2.5cm}

% === Moderne Schriftarten ===
\setmainfont{TeX Gyre Termes}[          % Times-ähnlich, frei verfügbar
  Ligatures=TeX,
  Numbers=OldStyle
]
\setsansfont{TeX Gyre Heros}[           % Helvetica-ähnlich
  Ligatures=TeX,
  Scale=MatchLowercase
]
\setmonofont{TeX Gyre Cursor}[          % Courier-ähnlich
  Scale=MatchLowercase
]

% === Hyperref-Konfiguration ===
\hypersetup{
  unicode=true,
  colorlinks=true,
  linkcolor=blue!70!black,
  urlcolor=blue!80!black,
  citecolor=green!60!black,
  pdftitle={Projektstruktur der modernisierten LaTeX-Vorlage},
  pdfauthor={LaTeX Template Documentation},
  pdfsubject={LuaLaTeX Template Structure},
  pdfkeywords={LuaLaTeX, Template, Struktur, Modernisierung}
}

% === Moderne Farben ===
\usepackage{xcolor}
\definecolor{codeblue}{RGB}{0, 102, 153}
\definecolor{success}{RGB}{40, 167, 69}
\definecolor{warning}{RGB}{255, 193, 7}

\title{Projektstruktur der LaTeX-Vorlage\\
{\large LuaLaTeX-optimiert für 2026}}
\author{}
\date{\today}

\begin{document}

\maketitle

\section{Wurzelverzeichnis}

\textcolor{codeblue}{\texttt{1Main.tex}}
\begin{itemize}
    \item Hauptdokument
    \item Bindet alle Kapitel, Verzeichnisse und Styles ein
    \item \textbf{Wichtig:} Kompilierung mit \texttt{lualatex --shell-escape Main.tex}
    \item \textcolor{success}{✓} Vollständig LuaLaTeX-kompatibel
\end{itemize}

\section{\texttt{content/} – Die Kapitel und Inhalte}

Hier wird der Hauptinhalt der Arbeit geschrieben. Die wichtigsten Dateien:

\begin{longtblr}{
  colspec = {X[l] X[l]},
  width = \textwidth,
  hlines, vlines,
  row{1} = {font=\bfseries, bg=gray!20}
}
Datei & Inhalt \\
\texttt{Einleitung.tex} & Template mit Beispiel-Einleitung \\
\texttt{Grundlagen.tex} & Template mit theoretischen Grundlagen \\
\texttt{Hauptteil.tex} & Template für Implementierung/Konzept \\
\texttt{Schluss.tex} & Template für Fazit und Ausblick \\
\texttt{Code.tex} & Code-Beispiele (Anhang) \\
\texttt{allgemein/Abstract.tex} & Abstract (optional) \\
\texttt{allgemein/Aufgabenstellung.tex} & Aufgabenstellung \\
\texttt{allgemein/Danksagung.tex} & Danksagung (optional) \\
\texttt{allgemein/Glossar.tex} & Glossar-Definitionen \\
\end{longtblr}

\section{\texttt{default\_frame/} – Verzeichnisse und formale Abschnitte}

Diese Dateien regeln die formale Struktur der Arbeit:

\begin{longtblr}{
  colspec = {X[l] X[l]},
  width = \textwidth,
  hlines, vlines,
  row{1} = {font=\bfseries, bg=gray!20}
}
Datei & Zweck \\
\texttt{Titelseite.tex} & Layout der Titelseite\\
\texttt{Inhaltsverzeichnis.tex} & Inhaltsverzeichnis\\
\texttt{Abbildungsverzeichnis.tex} & Liste aller Abbildungen\\
\texttt{Tabellenverzeichnis.tex} & Liste aller Tabellen\\
\texttt{Abkuerzungsverzeichnis.tex} & Glossaries-basiert\\
\texttt{Literaturverzeichnis.tex} & Biblatex-basiert\\
\texttt{Abbildungsquellenverzeichnis.tex} & Bildquellen (biblatex-Kategorien) \\
\texttt{Glossar.tex} & Fachbegriffe-Verzeichnis\\
\texttt{SelbststaendigkeitserklaerungUniBwM.tex} & Rechtliche Erklärung\\
\end{longtblr}

\section{\texttt{default\_style/} – Stildefinitionen}

\begin{longtblr}{
  colspec = {X[l] X[l]},
  width = \textwidth,
  hlines, vlines,
  row{1} = {font=\bfseries, bg=gray!20}
}
Datei & Modernisierung \\
\texttt{load\_all.tex} & ✓ PDF-Befehle entfernt, bereinigt \\
\texttt{documentclass.tex} & ✓ pdftex-Option entfernt, 12pt, LuaTeX-Erkennung \\
\texttt{Schriftart-settings.tex} & ✓ fontspec + Unicode-Fonts + polyglossia \\
\texttt{load\_packages.tex} & ✓ Kompletter Umbau: unicode-math, tabularray, etc. \\
\texttt{Seitenrand-settings.tex} & ✓ Responsive Modi (Druck/Digital) \\
\texttt{Kopf\_und\_Fusszeile\_misc.tex} & ✓ Farbintegration, moderne Stile \\
\texttt{Farben\_definieren-settings.tex} & ✓ Corporate Design, semantische Farben \\
\texttt{Listings-settings.tex} & ✓ Komplett ersetzt durch minted (300+ Sprachen!) \\
\texttt{Hyperref-settings.tex} & ✓ LuaLaTeX-Metadaten, farbige Links \\
\texttt{Abkuerzungsverzeichnis-settings.tex} & ✓ Glossaries-System modernisiert \\
\texttt{TocStyle.tex} & ✓ KOMA-optimiert, Farbintegration \\
\texttt{Fussnoten-settings.tex} & ✓ Moderne Formatierung, Farbunterstützung \\
\texttt{Absatz\_Abstand-settings.tex} & ✓ KOMA-parskip, responsive Modi \\
\texttt{Emoji-settings.tex} & ✓ NEU: Unicode-Emojis + semantische Symbole \\
\end{longtblr}

\subsection{Reparierte .sty-Dateien}

\begin{longtblr}{
  colspec = {X[l] X[l]},
  width = \textwidth,
  hlines, vlines,
  row{1} = {font=\bfseries, bg=gray!20}
}
Datei & Status \\
\texttt{type1ec.sty} & ✓ LuaLaTeX-kompatible Version (fontspec-basiert) \\
\texttt{bold-extra.sty} & ✓ FakeBold für moderne Unicode-Fonts \\
\texttt{picins.sty} & ✓ Ersetzt durch wrapfig + tcolorbox \\
\end{longtblr}

\section{\texttt{bibliography/} – Literaturverwaltung}

\begin{longtblr}{
  colspec = {X[l] X[l]},
  width = \textwidth,
  hlines, vlines,
  row{1} = {font=\bfseries, bg=gray!20}
}
Datei & Verwendung \\
\texttt{references.bib} & Moderne biblatex-Datenbank \\
& Kategorien für Bildquellen unterstützt \\
& Unicode-Zeichen funktionieren \\
\end{longtblr}

\section{Neue Features des modernisierten Systems}

\subsection{Code-Highlighting mit minted}
\begin{itemize}
    \item ✓ 300+ Programmiersprachen (Rust, Dart, Scala, etc.)
    \item ✓ Perfektes Syntax-Highlighting (Pygments-basiert)
    \item ✓ Unicode-Support (Emojis in Code funktionieren)
    \item ✓ Automatisches Caching für schnelle Kompilierung
\end{itemize}

\subsection{Emoji-System}
\begin{itemize}
    \item ✓ Native Unicode-Emojis: �� �� ��
    \item ✓ Semantische Befehle: \texttt{\textbackslash success}, \texttt{\textbackslash warning}
    \item ✓ Automatische Font-Erkennung
\end{itemize}

\subsection{Responsive Design}
\begin{itemize}
    \item ✓ Automatische Druck/Digital-Modi
    \item ✓ Intelligente Seitenränder
    \item ✓ Farbschema-Integration
\end{itemize}

\section{Kompilierung}

\subsection{Neue Kommandos}
\begin{longtblr}{
  colspec = {X[l] X[l]},
  width = \textwidth,
  hlines, vlines,
  row{1} = {font=\bfseries, bg=gray!20}
}
Alt (pdfLaTeX) & Neu (LuaLaTeX) \\
\texttt{pdflatex Main.tex} & \texttt{lualatex --shell-escape Main.tex} \\
\texttt{makeindex} & \texttt{makeglossaries Main} \\
Keine Emoji-Fonts nötig & Python + Pygments erforderlich \\
\end{longtblr}

\subsection{Vollständige Kompilierung}
\begin{enumerate}
    \item \texttt{lualatex --shell-escape Main.tex}
    \item \texttt{makeglossaries Main}
    \item \texttt{lualatex --shell-escape Main.tex}
    \item \texttt{lualatex --shell-escape Main.tex}
\end{enumerate}

\section*{Zusammenfassung: Moderne LaTeX-Vorlage 2026}

\begin{longtblr}{
  colspec = {X[l] X[l] X[l]},
  width = \textwidth,
  hlines, vlines,
  row{1} = {font=\bfseries, bg=gray!20}
}
Bereich & Neu (2026) \\
Engine & LuaLaTeX \\
Fonts & fontspec + Unicode \\
Code & minted (300+ Sprachen) \\
Sprache & polyglossia \\
Tabellen & tabularray \\
Farben & xcolor + Corporate Design \\
Emojis & Native Unicode \\
Modernität & 3/5 & 5/5 \\
\end{longtblr}

\vspace{1cm}

\textcolor{success}{\textbf{✓ Diese Vorlage ist vollständig zukunftssicher und einsatzbereit!}}

\begin{center}
\textit{Erstellt als moderne Template-Vorlage 2026}\\
\textit{Basierend auf LuaLaTeX mit modernsten Paketen}\\
\textit{Für Bachelor-, Master- und Doktorarbeiten an der UniBw München}
\end{center}

\end{document}